% Interior study for a collection of short pieces by Muhib.
% Editorial titles and order are provisional; source map: PREVIEW_NOTES.md.
\documentclass[11pt,twoside,openany]{book}
\usepackage[T1]{fontenc}
\usepackage[utf8]{inputenc}
\usepackage{mathpazo}
\usepackage{microtype}
\usepackage{geometry}
\usepackage{fancyhdr}
\usepackage{needspace}
\geometry{paperwidth=5.5in,paperheight=8.5in,inner=0.82in,outer=0.68in,top=0.82in,bottom=0.78in,headheight=13.6pt,headsep=0.27in}
\linespread{1.12}
\setlength{\parindent}{0pt}
\setlength{\parskip}{0.48\baselineskip}
\widowpenalty=10000
\clubpenalty=10000
\fancyhf{}
\fancyhead[LE]{\small\scshape Muhib}
\fancyhead[RO]{\small\scshape A Handful of Days}
\fancyfoot[LE,RO]{\thepage}
\renewcommand{\headrulewidth}{0pt}
\pagestyle{fancy}
\fancypagestyle{plain}{\fancyhf{}\fancyfoot[C]{\thepage}\renewcommand{\headrulewidth}{0pt}}
\newcommand{\piecehead}[1]{%
  \Needspace{6\baselineskip}%
  \par\addvspace{1.2\baselineskip}%
  {\noindent\large\bfseries #1\par}%
  \vspace{-0.15\baselineskip}%
}
\begin{document}
\thispagestyle{plain}
\vspace*{0.42in}
{\noindent\small\scshape A selection of short pieces\par}
\vspace{0.4in}
{\noindent\Huge\bfseries A Handful of Days\par}
\vspace{0.2in}
\noindent\rule{0.75in}{0.4pt}
\vspace{0.26in}

% CVThrYVI3Pu
\piecehead{The Cinema, Alone}
I went to the cinema alone today and was forced to internalise my experience before sharing it with anyone else. It was refreshing. We live in a share culture; we share our thoughts, feelings and even our heart. It's nice to pause that for a while and share with the self first.

% Cw4pSf5oSDC
\piecehead{Thirty}
I remember once complaining about my age. I was maybe seven, and I proclaimed that I wished I was thirty. The adults at the dinner table burst into laughter, practically entering into hysterics. They urged me to repeat it, and then cackled together, huddled over in amusement. At the time, I had no clue; I just assumed they found it hilarious. I thought, wow, thirty must be amazing, look how happy it makes them.

Now that I'm thirty, I can confirm they were not laughing out of joy. They were laughing at innocence, all harboring a secret desire to be seven again.

% Coc6msLoJPS
\piecehead{The Open Door}
I remember quite vividly the time my father told me about his youth, when he used to walk into his neighbor's house. I was shocked and asked, ``Did you knock on the door?'' He replied, ``No, I didn't need to. The door was always open.'' I asked, ``Why did you go in?'' He said, ``We needed sugar, so I took some from there.'' I stood there in disbelief, scratching my head and looking left and right.

I said with great confusion, ``Did you ask for permission first?'' He laughed and replied, ``No, I didn't need to.'' He looked at me and smiled, and I smiled back. He then said, ``Those times don't exist anymore. We lock our doors and draw our curtains, but back then, your house was my house. We were a family, and we all felt safe.''

% CZJbdAVowH7
\piecehead{Every Party}
Death has been invited to every single party of life.

% CZRWQXuIOiK
\piecehead{What Is Missing}
Through the lens of capitalism you are always missing something. You lack, and for that you must torture yourself until you figure out what it is you lack. Conveniently, what you lack is being sold to you, and with the wages you earned from torturing yourself you can buy it.

% DdNW9cQiCYh
\piecehead{Grace}
Hustling is exhausting. Forcing is exhausting. Pushing that which is not blessed to move---exhausting. Over the last few years, I have been almost obsessed with the idea of grace. The phenomenal truth is that even the smallest things, once touched with grace, can have a profound impact. In this world of metrics, views, likes, and comments, we often mistake the spectacle for success---but real success is found in the quiet moments of grace, where what we have goes a long way. In a world consumed by numbers, production, and quantity, there is also a spiritual world concerned with quality. God can bless us, and nothing can stop that blessing. No structure can limit it---perhaps that's why it's called grace.

% CtaEaapoBYE
\piecehead{The First Prayer}
I was probably 16 when I first tried to pray. I had no idea what I was doing. It was confusing, embarrassing, but also strangely familiar. Want to know the funny part? I was praying wrong for almost a year. Everything was mixed up. And guess what else? I was facing the wrong direction. I got so upset and almost gave up, but then someone close to me said, ``It's okay to be a beginner.'' So, if you're a beginner like me, remember, it's okay. You don't need to be perfect. With time, we'll improve.

% CcKUq5Ioa7G
\piecehead{A Door}
A few years ago I told a wise spiritual man that the door to God's vast and beautiful mercy was open. His mouth formed a subtle curve to the right. I was a little confused. I asked him, ``Why are you smiling?'' He turned and said, ``What makes you think there's a door?''

% CmKFe-PIdQa
\piecehead{Bread in the Mountains}
I remember seeing a distressed elderly lady in Morocco. She lived in the mountains. I don't know where. But it was remote. She was crying. She held her face in her hands. We spotted her from further up the road and decided to give her some bread that we happened to be gifted earlier on. Bread in Morocco is a huge circular loaf. It is also seen as a spiritual manifestation of God's immense grace.

When she saw that we had given it to her she started crying even more. We stood there a little confused so I decided to ask what was wrong. She told us that her son had just died. She told us that only he looked after her. She said that she had no money for her electricity bill. She looked up to the sky and said she had been begging God for help and that this bread was a sign.

It was a sign that God had not deserted her, and that she would be okay. One of the brothers with me decided to pay an entire month's electricity bill for her. That is when things got really intense. She sat there in a spiritual ecstasy and said, ``Oh my Lord! How wonderful you are! Here I am asking you, begging you, seeking your help and here you have sent me Your support. My sustenance was destined to be delivered by strangers who live many many miles away from me! How can we despair? How can we despair when you have indeed promised that you are generous!''

She said that the entire moment was a miracle and an Ayah. Ayah here means a sign, a sign that God is indeed what He promises himself to be---Al Kareem; The Most Generous. It is a moment that will forever continue to shock and amaze me.

She reminded us that everything would ultimately be okay. She was short with rosy cheeks and looked like the grandmother you would look to for comfort.

% C1dtfCboJnS
\piecehead{More Than Content}
Palestine is not content. Palestinians are more than images and videos. Social media is neither your friend nor mine. This platform should not deceive you into believing that your activism ends after you have viewed and shared the daily show. This is real life, both horrific and atrocious. The devastation in Gaza cannot be fully captured by any phone screen or camera. Instagram and tech companies cannot fool you into believing that activism relies on algorithms. It doesn't. Our activism, support, and desire for the dignity of Palestinian lives depend on our intentions, knowledge, and right action. Don't let them fool you.

% CqMo8uwolct
\piecehead{I Don't Know}
There's a real beauty in raising your hands and saying, ``I don't know.'' The start of all wisdom is admitting we don't know. In today's world, there is so much pressure to always know, to always be right, and to always have the answers. The truth is, many times, we don't have the answers. We don't know. When we admit it, that's when we create the chance to learn.

% C17kJLxIYiV
\piecehead{Baraka}
I had a strange chest feeling when I learned about `baraka'. It's still a mystery, but I can tell you what I was told. It was simple. Direct. Simply put, ``Baraka makes that which you have last; it expands the small; it brings blessing to each part.'' The end. My worldview changed drastically. Yes, I sometimes forget. But the idea that even a tiny thing can be blessed is revolutionary. Now, having enough is about God's grace, not quantity. Maybe the feeling in my chest was relief---relief that God can bless less than much. Relief that perhaps the anxiety of hustling and forcing can melt away, and perhaps, Baraka can take over and multiply what I cannot.

% Ce_wA01o2Xo
\piecehead{Together}
I am guessing it is safe to say we're all trying to hold hope. We're all riding, betting, leaning on hope. It's what brings us together. It's why we try to smile after sighing. We may never say it loud to each other, but we are all holding onto this huge rope of hope.

Together.

\end{document}
